% Balanced complete-edition layout; entries unchanged from the accepted bibliography.
\begin{thebibliography}{99}
\label{refs:chapter02}
\bibitem[Benkler(2002)]{benkler2002}Benkler, Y. (2002).
Coase's penguin, or, Linux and the nature of the firm.
\emph{Yale Law Journal}, 112, 369--446.
\url{https://yalelawjournal.org/article/coases-penguin-or-linux-and-the-nature-of-the-firm}.

\bibitem[Borgatti(2005)]{borgatti2005}Borgatti, S.~P. (2005).
Centrality and network flow.
\emph{Social Networks}, 27(1), 55--71.
\doi{10.1016/j.socnet.2004.11.008}.

\bibitem[Bowman and Ambrosini(2000)]{bowman2000}Bowman, C., and Ambrosini, V. (2000).
Value creation versus value capture: Towards a coherent definition of value in strategy.
\emph{British Journal of Management}, 11(1), 1--15.
\doi{10.1111/1467-8551.00147}.

\bibitem[Coelho and Valente(2017)]{coelho2017}Coelho, J., and Valente, M.~T. (2017).
Why modern open source projects fail.
In \emph{Proceedings of the 2017 11th Joint Meeting on Foundations of Software Engineering}, pp.~186--196. ACM.
\doi{10.1145/3106237.3106246}.

\bibitem[Dabbish et~al.(2012)]{dabbish2012}Dabbish, L., Stuart, C., Tsay, J., and Herbsleb, J. (2012).
Social coding in GitHub: Transparency and collaboration in an open software repository.
In \emph{Proceedings of the ACM 2012 Conference on Computer Supported Cooperative Work}, pp.~1277--1286. ACM.
\doi{10.1145/2145204.2145396}.

\bibitem[Decan et~al.(2019)]{decan2019}Decan, A., Mens, T., and Grosjean, P. (2019).
An empirical comparison of dependency network evolution in seven software packaging ecosystems.
\emph{Empirical Software Engineering}, 24(1), 381--416.
\doi{10.1007/s10664-017-9589-y}.

\bibitem[Espeland and Sauder(2007)]{espeland2007}Espeland, W.~N., and Sauder, M. (2007).
Rankings and reactivity: How public measures recreate social worlds.
\emph{American Journal of Sociology}, 113(1), 1--40.
\doi{10.1086/517897}.

\bibitem[Fang and Neufeld(2009)]{fang2009}Fang, Y., and Neufeld, D. (2009).
Understanding sustained participation in open source software projects.
\emph{Journal of Management Information Systems}, 25(4), 9--50.
\doi{10.2753/MIS0742-1222250401}.

\bibitem[Gousios et~al.(2008)]{gousios2008}Gousios, G., Kalliamvakou, E., and Spinellis, D. (2008).
Measuring developer contribution from software repository data.
In \emph{Proceedings of the 2008 International Working Conference on Mining Software Repositories}, pp.~129--132. ACM.
\doi{10.1145/1370750.1370781}.

\bibitem[Han et~al.(2024)]{han2024ose}Han, F., Zhao, S., Wang, W., Zhou, A., Qian, W., Zhou, X., Peng, J., You, L., Chen, Y., Xia, X., Tang, Y., Yang, L., and Tian, C. (2024).
Open Source Evaluatology: An evaluation framework and methodology for open source ecosystems based on evaluatology.
\emph{BenchCouncil Transactions on Benchmarks, Standards and Evaluations}, 4(4), 100190.
\doi{10.1016/j.tbench.2025.100190}.
\emph{Bibliographic note: the journal assigns this article to the December 2024 issue; the DOI and online publication date are in 2025.}

\bibitem[Han et~al.(2026a)]{han2026rank}Han, F., Zhao, S., Wang, W., Peng, J., and Xia, X. (2026a).
OpenRank: A centrality algorithm for high-dimensional heterogeneous networks in open source collaboration.
\emph{Information Sciences}, 730, 122893.
\doi{10.1016/j.ins.2025.122893}.

\bibitem[Howison and Crowston(2014)]{howison2014}Howison, J., and Crowston, K. (2014).
Collaboration through open superposition: A theory of the open source way.
\emph{MIS Quarterly}, 38(1), 29--50.
\doi{10.25300/MISQ/2014/38.1.02}.

\bibitem[Joblin et~al.(2017)]{joblin2017}Joblin, M., Apel, S., Hunsen, C., and Mauerer, W. (2017).
Classifying developers into core and peripheral: An empirical study on count and network metrics.
In \emph{Proceedings of the 2017 IEEE/ACM 39th International Conference on Software Engineering}, pp.~164--174. IEEE.
\doi{10.1109/ICSE.2017.23}.

\end{thebibliography}
\clearpage
\begingroup
\renewcommand{\bibsection}{\section*{References\enspace{\normalfont\sffamily\fontsize{11}{14}\selectfont continued}}}
\begin{thebibliography}{99}
\bibitem[Messick(1995)]{messick1995}Messick, S. (1995).
Validity of psychological assessment: Validation of inferences from persons' responses and performances as scientific inquiry into score meaning.
\emph{American Psychologist}, 50(9), 741--749.
\doi{10.1037/0003-066X.50.9.741}.

\bibitem[Shapley(1953)]{shapley1953}Shapley, L.~S. (1953).
A value for n-person games.
In H.~W. Kuhn and A.~W. Tucker (eds.), \emph{Contributions to the Theory of Games}, Vol.~II, pp.~307--318. Princeton University Press.
\doi{10.1515/9781400881970-018}.

\bibitem[von Hippel and von Krogh(2003)]{vonhippel2003}von Hippel, E., and von Krogh, G. (2003).
Open source software and the “private-collective” innovation model: Issues for organization science.
\emph{Organization Science}, 14(2), 209--223.
\doi{10.1287/orsc.14.2.209.14992}.

\bibitem[von Krogh et~al.(2012)]{vonkrogh2012}von Krogh, G., Haefliger, S., Spaeth, S., and Wallin, M.~W. (2012).
Carrots and rainbows: Motivation and social practice in open source software development.
\emph{MIS Quarterly}, 36(2), 649--676.
\url{https://aisel.aisnet.org/misq/vol36/iss2/16/}.

\bibitem[Wang et~al.(2026)]{wang2026digger}Wang, W., Han, F., Zhao, S., Zhou, X., Qian, W., Zhou, A., Xia, X., Yang, L., Wang, R., Jiang, N., and Duan, M. (2026).
OpenDigger: A practical framework for assessing community health and sustainability in open source collaboration platforms.
In \emph{Proceedings of the ACM Web Conference 2026}, pp. 7745--7756. ACM.
\doi{10.1145/3774904.3792807}.

\bibitem[Xia et~al.(2022)]{xia2022}Xia, X., Weng, Z., Wang, W., and Zhao, S. (2022).
Exploring activity and contributors on GitHub: Who, what, when, and where.
In \emph{2022 29th Asia-Pacific Software Engineering Conference}, pp. 11--20. IEEE.
\doi{10.1109/APSEC57359.2022.00013}.

\bibitem[Xia(2024)]{xia2024thesis}Xia, X. (2024).
\emph{Data-Driven Empirical Research on the Key Elements of Sustainability in Open Source Software Communities}.
Doctoral dissertation, East China Normal University, March 2024. In Chinese; English title supplied by the dissertation.

\bibitem[Young et~al.(2021)]{young2021}Young, J.-G., Casari, A., McLaughlin, K., Trujillo, M.~Z., Hébert-Dufresne, L., and Bagrow, J.~P. (2021).
Which contributions count? Analysis of attribution in open source.
In \emph{Proceedings of the 18th IEEE/ACM International Conference on Mining Software Repositories}, pp.~242--253.
\doi{10.1109/MSR52588.2021.00036}.

\bibitem[Zhan et~al.(2026)]{zhan2026book}Zhan, J., Wang, L., Gao, W., Wang, C., Fan, F., Kang, G., and Li, H. (2026).
\emph{Evaluatology: The Science of Revealing Causes and Effects} [title translated from Chinese].
Chinese edition, version 1.0, August 2026. BenchCouncil Press.
ISBN 978-988-71596-9-8. \doi{10.66834/33wi1xx8}.
\emph{Page references refer to the printed pagination of this edition.}

\bibitem[Zhao et~al.(2024)]{zhao2024}Zhao, S., Xia, X., Fitzgerald, B., Li, X., Lenarduzzi, V., Taibi, D., Wang, R., Wang, W., and Tian, C. (2024).
OpenRank Leaderboard: Motivating open source collaborations through social network evaluation in Alibaba.
In \emph{Proceedings of the 46th International Conference on Software Engineering: Software Engineering in Practice}, pp. 346--357. ACM.
\doi{10.1145/3639477.3639734}.

\bibitem[Zhao(2025)]{zhao2025thesis}Zhao, S. (2025).
\emph{Research and Application of Key Technologies for Open Source Ecosystem Data Platform}.
Doctoral dissertation, Tongji University, September 2025. In Chinese; English title supplied by the dissertation.
\end{thebibliography}
\endgroup
